\section{Simulating DESI-like Images and Spectra}
\label{sec:data-generate}

\subsection{Galaxy and Observation Model}
\label{ssec:kl-model}

We simulate our galaxies following the DESI-KL forecast setup of \citep{KLDESI}, in which each galaxy is observed with a photometric image and a set of DESI fiber spectra. DESI takes spectra by placing optical fibers of 1.5$\arcsec$ diameter on its targets \citep{DESIOverview}. For KL, we need multiple fibers placed along a disk galaxy's major and minor axes in order to constrain its velocity field.

A similar fiber arrangement has already been used by the DESI Peculiar Velocity Survey (DESI-PV), which places offset fibers along the major axes of disk galaxies to measure their rotation velocities \citep{DESIPVTargetSel2023, DESIPVTF2025}. This makes spectra from DESI-PV suitable for a pilot KL measurement with DESI, which is being carried out in Willett et al. (in prep.). The method presented in this work could eventually be applied to such data as well.

\begin{table}
    \centering
    \begin{tabular}{llcl}
        \toprule
        Parameter & Description & Training prior & Units\\
        \tableline
        $g_1$ & Reduced shear & $\mathcal{U}$(-0.1, 0.1) & 1\\
        $g_2$ & Reduced shear & $\mathcal{U}$(-0.1, 0.1) & 1\\
        $\theta_\mathrm{int}$ & Position angle & $\mathcal{U}$(-$\pi$, $\pi$) & rad\\
        $\cos i$ & Inclination & $\mathcal{U}$(0, 1) & 1\\
        $v_0$ & Systemic velocity & $\mathcal{U}$(-30, 30) & km s$^{-1}$\\
        $\vcirc$ & Max. rotation velocity & $\mathcal{U}$(60, 540) & km s$^{-1}$\\
        $R_\mathrm{vscale}$ & Velocity scale radius & $\mathcal{U}$(0.1, 5.0) & arcsec\\
        $R_h$ & Half-light radius & $\mathcal{U}$(0.1, 5.0) & arcsec\\
        $\log_{10}F_{\mathrm{H}\alpha}$ & Central H$\alpha$ flux & $\mathcal{U}$(-17, -14) & \shortstack[l]{erg s$^{-1}$\\ cm$^{-2}$}\\
        \tableline
        $q_z$ & Disk aspect ratio & 0.1 & 1 \\
        $\sigma_\lambda^\mathrm{int}$ & Rest-frame line width & 0.05 & nm \\
        $z$ & Redshift & 0.3 & 1 \\
        \tableline
        \end{tabular}
    \caption{List of parameters used in our KL model. The top section lists the nine parameters inferred by our model, along with the priors from which the training and validation sets are drawn (Sec.~\ref{ssec:datasets}). $\mathcal{U}$(min, max) denotes a uniform distribution. $F_{\mathrm{H}\alpha}$ is the H$\alpha$ flux within the central fiber. The bottom section lists fixed parameters.}
    \label{tab:kl-params}
\end{table}

For our simulations, we use the KL forward-modeling pipeline of \citep{KLInferenceI}, which was modified in \citep{KLDESI} for DESI and Legacy Survey observations. Given a set of galaxy parameters (Table~\ref{tab:kl-params}), the simulator outputs a noiseless $r$-band image and five fiber spectra. We use \textsc{GalSim}~\citep{galsim} to generate an inclined exponential profile for the galaxy's intensity, specified by its half-light radius $R_h$, inclination $i$, and position angle $\theta_\mathrm{int}$, with a fixed aspect ratio of $q_z=0.1$. The profile is then sheared by $(g_1, g_2)$. For the velocity field, we follow \citep{KLDESI} and assume an infinitely thin circular disk with an arctan rotation curve \citep{Courteau97},
\begin{equation}
\label{eqn:vrot}
    v_\mathrm{rot}(R) = \frac{2\vcirc}{\pi}\tan^{-1}\left(\frac{R}{R_\mathrm{vscale}}\right),
\end{equation}
where $R$ is the distance from the center of the disk and $R_\mathrm{vscale}$ is the velocity scale radius. The LoS velocity field is then calculated from $v_\mathrm{rot}$, $i$ and $\theta_\mathrm{int}$, shifted by the systemic velocity $v_0$, and sheared in the same way as the intensity profile, following Eq.~(\ref{eqn:shear-vfield}).

For the spectral energy distribution (SED), we use a template galaxy continuum as in \citep{KLDESI}, scaled so that the galaxy's $r$-band AB magnitude is $m_r$, before adding a gaussian H$\alpha$ emission line on top. The H$\alpha$ flux $F_{\mathrm{H}\alpha}$ is left as a free parameter, which we define as the line flux that falls within the central fiber. As in \citep{KLDESI}, we fix the rest-frame line width to 0.05 nm and place all galaxies at a redshift of $z=0.3$, which puts H$\alpha$ in the $z$ band of the DESI spectrograph. We only simulate a single redshift. For an observed galaxy, we assume that its spectra would be cut out in a window centered on where the redshifted H$\alpha$ line should be, so the line always sits at the same position in the spectra. In Appendix \ref{app:redshift}, we verify that the shear estimates from our model remain unbiased for datavectors simulated at other redshifts, even though the kinematic parameters and H$\alpha$ flux are measured as if the galaxy were at $z=0.3$.

Using the intensity profile, velocity field, and SED, the simulator builds a 3D data cube with the 2D image position and wavelength, from which it generates the image and spectra. For a detailed description of this process, refer to \citep{KLDESI}. Our image is a $48\times48$ $r$-band stamp with the Legacy Survey pixel scale of 0.2637$\arcsec$, using the DECam $r$-band throughput and an airy PSF with 1.0$\arcsec$ FWHM. Each fiber spectrum covers a $\sim$5 nm window around the H$\alpha$ line at the DESI dispersion of 0.08 nm per pixel, giving us a $5\times64$ array for the five fibers of each galaxy. The instrument properties and observing configurations follow \citep{KLDESI}, and are listed in Table~\ref{tab:obs-conf}. We generate all images and spectra without noise, and add noise to them during training and inference (Sec.~\ref{ssec:datasets}).

\begin{table}
    \centering
    \begin{tabular}{lccl}
        \toprule
        Property & Image & Fiber & Units\\
        \tableline
        Telescope & CTIO Blanco & KPNO Mayall & - \\
        Effective diameter & 378.29 & 332.42 & cm \\
        Pixel scale & 0.2637 & - & arcsec\\
        PSF FWHM & 1.0 & 1.0 & arcsec \\
        Gain & 4 & 1 & e$^{-}$/ADU \\
        Dispersion & - & 0.08 & nm/pixel \\
        Exposure time & 60 & 180/600 & s \\
        Sky level & 44.54 & - & e$^{-}$ s$^{-1}$ pixel$^{-1}$ \\
        Read noise & 2.6 & - & e$^{-}$ \\
        Fiber radius & - & 0.75 & arcsec \\
        Fiber offset & - & 1.5 & arcsec \\
        Data dimensions & $48\times48$ & $5\times64$ & pixel \\
        \tableline
        \end{tabular}
    \caption{Instrument properties and observing configurations used to simulate our images and spectra. Fields with ``-'' are not applicable. Fiber exposure times are listed for the central/offset fibers respectively. The image sky level and read noise are only used for the targeting image that sets the fiber positions (Sec.~\ref{ssec:kl-model}).}
    \label{tab:obs-conf}
\end{table}

For each galaxy, we place one fiber on its center, two on either side of its observed major axis, and two on either side of its observed minor axis, all with a 1.5$\arcsec$ offset from the center. Following \citep{KLDESI}, the central fiber has an exposure time of 180 s, and each offset fiber has an exposure time of 600 s. While the true observed axes of the galaxy is directly computable from its shear, inclination and intrinsic position angle, we do not want to use those for our offset fiber placement, as that could leak more information than intended on the shear direction, especially for face-on galaxies \XW{Maybe an appendix to explain this}. Instead, we simulate realistic fiber placement by measuring the position angle on a noisy realization of each galaxy. In DESI-PV, the placement of offset fibers are determined by the position angle of the galaxy in the Siena Galaxy Atlas (SGA) \citep{DESIPVTargetSel2023}, which is measured from the second moments of the galaxy's light distribution in the Legacy Survey images \citep{SGA2020}. We emulate that process here by using \textsc{GalSim}'s HSM module to measure the adaptive second moments of a noised realization of each galaxy, with the noise level determined by the Legacy Survey $r$-band observing configuration of \citep{KLDESI} (Table~\ref{tab:obs-conf}).

\subsection{Training and Validation Data}
\label{ssec:datasets}

We generate 1 million galaxies for our training set and 100 thousand galaxies for our validation set. For both sets, we sample the nine parameters in Table~\ref{tab:kl-params} from uniform distributions, using Latin Hypercube sampling to achieve even sampling across parameter space. Contrary to \citep{KLDESI}, we sample the inclination in $\cos i$ instead of $\sin i$, since on the sky we see randomly oriented disks, which is uniformly distributed in $\cos i$. We sample the H$\alpha$ flux uniformly in $\log_{10}F_{\mathrm{H}\alpha}$. The $r$-band magnitude $m_r$, which sets the brightness of the continuum, is sampled separately from $\mathcal{U}(15, 23.4)$, where the upper limit corresponds to the median $5\sigma$ depth of the Legacy Survey $r$-band images for an extended source \citep{KLDESI}. The parameter ranges are chosen to be wide enough to contain the vast majority of galaxies in our test sets (Sec.~\ref{ssec:test-set}). The data generation pipeline is highly parallelizable, using a total of $\sim$2000 CPU hours for the training set.

As mentioned in Sec.~\ref{ssec:kl-model}, we store our images and spectra in their noiseless form, and add noise to them on-the-fly during training and inference. This lets us draw new noise levels and realizations for every galaxy at every epoch, so that our model learns from a variety of noise levels for any set of parameters. Doing this also keeps our model as agnostic  as possible to observing conditions, since the noise level is specified directly by an independently drawn signal-to-noise ratio (SNR) rather than by exposure time or sky level. We use a Gaussian noise model for computational efficiency, which we believe is an adequate approximation for ground-based surveys for the scope of this work.

We draw the SNR for image and spectra independently. We define $\rho_I$ as the SNR of the image as a whole, where each pixel contributes according to how much of the galaxy's signal it contains:
\begin{equation}
\label{eqn:snr-image}
    \rho_I = \frac{\lVert \bm{M}_\mathrm{I} \rVert}{\sigma_I},
\end{equation}
where $\lVert \bm{M}_\mathrm{I} \rVert$ is the L2 norm of the noiseless image and $\sigma_I$ is the standard deviation of the noise in each pixel. Similarly, we define $\rho_S$ on the H$\alpha$ line in the central fiber,
\begin{equation}
\label{eqn:snr-spec}
    \rho_S = \frac{\lVert \bm{M}_{\mathrm{S},c} - \tilde{M}_{\mathrm{S},c} \rVert}{\sigma_S},
\end{equation}
where $\bm{M}_{\mathrm{S},c}$ is the noiseless central fiber spectrum, $\tilde{M}_{\mathrm{S},c}$ is its median flux, which we use as an estimate of the continuum level, and $\sigma_S$ is the noise standard deviation in each pixel of the central fiber. Since noise grows as the square root of exposure time, the noise in each offset fiber is scaled up by $\sqrt{600/180}$ relative to the central fiber. During training, we draw $\rho_I \sim \mathcal{U}(10, 1000)$ and $\rho_S \sim \mathcal{U}(1, 150)$ for each galaxy at every epoch. For the validation set, we keep the SNRs and noise realizations fixed across epochs, so that the validation loss can be most accurately compared between epochs.

In addition to the images and spectra, we also pass a set of metadata into our model: the positions of the five fibers on the image, $m_r$, $\rho_I$, and $\rho_S$. \XW{We pass the true $m_r$ and SNRs as metadata. In real data these would be measured with some uncertainty. Will be addressed.} Before passing our data into the model, we normalize each image by its L2 norm, and the five spectra of each galaxy by their combined L2 norm, which conserves the relative flux between fibers. The fiber positions are divided by the fiber offset of 1.5$\arcsec$, and $m_r$, $\rho_I$, and $\rho_S$ are linearly rescaled to $[-1, 1]$ using their training ranges. The purpose of this is to stabilize the training and keep our model robust against different observing configurations. As an example, different exposure times on the same galaxy would give different absolute pixel fluxes in the image, but should have no effect on the final shear measurement. Instead, the brightness of the galaxy and the quality of the data are available to the model through $m_r$ and the SNRs.

\subsection{DESI-like Test Sets}
\label{ssec:test-set}

In our final diagnostics, we would like to test our model on a galaxy population that matches as much as possible with the one that would be observed by a DESI-KL survey. To best compare with previous results, we build our test sets from the galaxy samples selected in \citep{KLDESI}, specifically samples 3 and 5. Both samples are selected from the DESI Bright Galaxy Survey (BGS) \citep{BGS}, with sample 3 containing all BGS galaxies (BGS Any) and sample 5 containing only BGS Bright galaxies. Additionally, both samples require a disk-like morphology in the Legacy Survey imaging \citep{LSImaging} and a half-light radius of $R_h \geq 1\arcsec$. Sample 5 is the more stringent of the two, keeping brighter galaxies but a lower number density. We will refer to the test sets built from these samples as Xu3 and Xu5.

We do not use sample 1 of \citep{KLDESI}, which selects Legacy Survey galaxies in general with $r \leq 22$. This is because our test sets require an H$\alpha$ SNR for each galaxy, which is only available for galaxies with a DESI spectrum, and as of DESI Data Release 1, most sample 1 galaxies do not have one. After requiring an H$\alpha$ measurement, the remaining sample 1 galaxies are very similar to those in sample 3. We also do not use samples 2, 4, and 6, which apply a stricter cut of $R_h \geq 1.5\arcsec$, as we did not see significant performance changes in our model between these two cuts.

For each sample, we compile a catalog of all galaxies that pass its selection cuts, and keep only the ones with an H$\alpha$ flux measurement from \textsc{FastSpecFit} \citep{FastSpecFit} in DESI Data Release 1 \citep{DESIDR1}. This catalog gives us empirical joint distributions of $m_r$, $R_h$, image SNR, and H$\alpha$ SNR (see Fig.~\ref{fig:test-set-samples}). We take the image SNR to be the SNR of the Legacy Survey $r$-band flux, and the H$\alpha$ SNR to be that of the \textsc{FastSpecFit} H$\alpha$ flux, and use them as our $\rho_I$ and $\rho_S$ respectively. We further restrict the catalog to the ranges covered by our training set: $15 \leq m_r \leq 23.4$, $0.1\arcsec \leq R_h \leq 5\arcsec$, $10 \leq \rho_I \leq 1000$, and $1 \leq \rho_S \leq 150$, which leaves 2.3 million galaxies in sample 3 and 1.7 million galaxies in sample 5.

From each catalog, we then randomly draw 100 thousand galaxies, taking $m_r$, $R_h$, $\rho_I$, and $\rho_S$ from the same galaxy so that the correlations between them are preserved. For each galaxy, we sample $\vcirc$ from the TFR given $m_r$ (Sec.~\ref{ssec:tf-prior}), so that the rotation velocities in our test sets follow a realistic relation with brightness. The rest of the KL parameters, $(g_1, g_2, \theta_\mathrm{int}, \cos i, v_0, R_\mathrm{vscale}, F_{\mathrm{H}\alpha})$, are drawn from the same priors as our training set (Table~\ref{tab:kl-params}). As with our training set, all test galaxies are simulated at $z=0.3$ with the observing configuration of Sec.~\ref{ssec:kl-model}, regardless of their catalog redshift. The resulting test sets have median values of $m_r=19.2$, $R_h=1.6\arcsec$, $\rho_I=174$, and $\rho_S=26$ for Xu3, and $m_r=19.0$, $R_h=1.8\arcsec$, $\rho_I=204$, and $\rho_S=27$ for Xu5.

\begin{figure*}
    \includegraphics[width=\linewidth]{figs/xu_samples_3_5_distributions.png}
    \caption{Redshift, $r$-band magnitude, half-light radius, and H$\alpha$ SNR distributions of samples 3 and 5 from \citep{KLDESI}, for galaxies with H$\alpha$ measurements in DESI DR1. We draw $m_r$, $R_h$, image SNR, and H$\alpha$ SNR jointly from these catalogs to build our Xu3 and Xu5 test sets. The catalog redshift is shown for reference only, as all test galaxies are simulated at $z=0.3$.}
    \label{fig:test-set-samples}
\end{figure*}
